\documentclass[10pt,a4paper]{amsart}
\usepackage[margin=19mm]{geometry}
\usepackage{amsmath,amssymb,mathtools}
\usepackage[scr=rsfs,cal=euler]{mathalfa}
\usepackage{xfrac}
\usepackage[unicode,hidelinks,bookmarks=false]{hyperref}
\newcommand{\ie}{i.e.\ }
\newcommand{\cf}{cf.\ }
\newcommand{\resp}{respectively\ }
\newcommand{\Subsection}{\subsection}
\providecommand{\nobreakdash}{\nobreak\hbox{-}}
\setlength{\emergencystretch}{2em}
\multlinegap0pt
\usepackage{xcolor}
\newtheorem{lem}{Lemma}
\newtheorem{prop}{Proposition}
\DeclareMathOperator{\osc}{{\rm osc}}
\title{A variation on the P\'{o}lya-Seg\H{o} principle in one dimension}
\author{Sorina Barza, Martin Lind}
\date{Source: arXiv:2607.03450v3 (CC BY 4.0)}
\begin{document}
\maketitle

\noindent\emph{Source and licence.} A variation on the P\'{o}lya-Seg\H{o} principle in one dimension. Version arXiv:2607.03450v3 (CC BY 4.0), distributed by arXiv under the Creative Commons Attribution 4.0 International licence, http://creativecommons.org/licenses/by/4.0/, as recorded in the arXiv licence metadata for this exact version and shown on the abstract page https://arxiv.org/abs/2607.03450v3. Full licence notice: \texttt{licenses/arxiv-2607.03450-CC-BY-4.0.txt}.\par\smallskip

\noindent\textit{Editorial scope.} Counted unit: the article's prop simplePS (source0.tex lines 341--353). The article's complete original proof of it is delivered with it (source0.tex lines 354--400, one \texttt{proof} environment of 47 lines). No mathematics is written, repaired or re-derived: the statement and the proof are the article's own lines, in the article's own notation, and no retained line is substituted or re-flowed.  Retained uncounted as the source-local context the proof needs: lines 147--159; lines 218--229.  Nothing was omitted from any retained span, so the package carries no omission record.  No retained line contains a citation, so the wrapper carries no bibliography and no bibliography entry.  No retained line contains a figure environment, an image-inclusion command or drawing code, so no image is delivered and the package declares no asset dependency.  Editorial formatting only, in addition: an AMS \texttt{amsart} preamble that restates the article's own declarations and macros used by the retained lines, namely the environments \texttt{lem}, \texttt{prop} on separate counters, together with the standard packages those lines need.  The retained environments print in this document as Lemma 1, Proposition 1 in order of appearance, whereas the article prints the same items with the numbering of its own full document.  The complete native source of the article as distributed in the arXiv e-print for this exact version is bundled unchanged as \texttt{sources/204386/source0.tex}.
\begin{lem}
    \label{decreasingLem}
    For any $y\ge0$
    \begin{equation}
        \label{eqlevelset}
        |\{x\in I:|f(x)|=y\}|=|\{t\in I^*:f^*(t)=y\}|.
    \end{equation}
    In particular, $f^*$ is strictly decreasing on $I^*$ if and only if 
    $$
    |\{x\in I:|f(x)|=y\}|=0
    $$
    for every $y\ge0$.
\end{lem}

\begin{lem}
    \label{mainLemma}
    Let $f:I\rightarrow\mathbb{R}$ be a continuous function. Take $y',y''\in f(I)$, $y'<y''$, and consider the set
    $$
    E=\{x\in I:y'<f(x)<y''\}.
    $$
    Then there exists an interval $J=(x',x'')\subset E$ such that
    \begin{equation}
    \label{extremal}
        \inf_{x\in J}f(x)=y'\quad\text{and}\quad\sup_{x\in J}f(x)=y''.
    \end{equation}
\end{lem}

\begin{prop}
    \label{simplePS}
    Let $1<p<\infty$ and $0<\alpha\le 1/p'$. Assume that $f\in C(I)$, $f\ge0$ on $I$, and 
    \begin{equation}
        \label{simplePSA3}
        |\{x\in I:f(x)=y\}|=0\quad\text{for every } y\ge0.
    \end{equation}
    Then 
    \begin{equation}
        \label{simplePS-ineq}
        \mathcal{V}_p^\alpha(f^*;I^*)\le\mathcal{V}_p^\alpha(f;I).
    \end{equation}
\end{prop}

\begin{proof}
    The assumption (\ref{simplePSA3}) implies that $f^*$ is strictly decreasing, by Lemma \ref{decreasingLem}. Further, $f$ and $f^*$ are equimeasurable. Therefore, for any $t\in I^*$
    \begin{equation}
        \label{simplePS-pf1}
        \mu_f(f^*(t))=\mu_{f^*}(f^*(t))=t.
    \end{equation}
    Take an arbitrary finite family of disjoint intervals $[a_k,b_k]\subset I^*~~(1\le k\le n)$ and consider
    $$
    S:=\sum_{k=1}^n\frac{(\osc(f^*;(a_k,b_k))^p}{(b_k-a_k)^{\alpha p}}=\sum_{k=1}^n\frac{(f^*(a_k)-f^*(b_k))^p}{(b_k-a_k)^{\alpha p}}.
    $$
     We shall find disjoint intervals $I_k~~(1\le k\le n)$ such that
    \begin{equation}
        \label{ineq1}
        \frac{(f^*(a_k)-f^*(b_k))^p}{(b_k-a_k)^{\alpha p}}\le\frac{\osc(f;I_k)^p}{|I_k|^{\alpha p}}\quad(1\le k\le n).
    \end{equation}
    Hence $S\le\mathcal{V}_p^\alpha(f;I)$ and  (\ref{simplePS-ineq}) follows.
    Fix $k\in\{1,2,\ldots,n\}$. Define
    \begin{equation}
        \nonumber
        M_k=\{x\in I: f^*(a_k)>f(x)>f^*(b_k)\}.
    \end{equation}
    Note that the sets $M_k$ are disjoint. Further, since 
    \begin{equation}
        \nonumber
        M_k=\{x\in I:f(x)> f^*(b_k)\}\setminus\{x\in I:f(x)\ge f^*(a_k)\}
    \end{equation}
    we have
    \begin{equation}
        \nonumber
        |M_k|=\mu_f(f^*(b_k))-\left(\mu_f(f^*(a_k))+|\{x\in I:f(x)=f^*(a_k)\}|\right)
    \end{equation}
    By (\ref{simplePSA3}), $|\{x\in I:f(x)=f^*(a_k)\}|=0$. Therefore, by (\ref{simplePS-pf1})
    $$
        |M_k|=\mu_f(f^*(b_k))-\mu_f(f^*(a_k))=b_k-a_k
    $$
    Since $f$ is continuous, we may invoke Lemma \ref{mainLemma} to obtain an open interval $I_k\subset M_k$ such that
    \begin{equation}
        \nonumber
        f^*(a_k)-f^*(b_k)={\rm osc}(f;I_k).
    \end{equation}
    Since $I_k\subset M_k$, there also holds $|I_k|\le b_k-a_k$, so
    \begin{equation}
        \nonumber
        \frac{(f^*(a_k)-f^*(b_k))^p}{(b_k-a_k)^{\alpha p}}\le\frac{\osc(f;I_k)^p}{|I_k|^{\alpha p}}.
    \end{equation}
    This concludes the proof.
\end{proof}

\end{document}
